\newpage


%----------------------------------------------------
\section{Training procedure} 
\label{ap:training_curves}

Each model is trained with the hyperparameters in \ref{ap:model_hyper}. Loss curves averaged over 5 runs are displayed for each experiment. Plots are with semi-log scales. 

\begin{figure}[H]
\centering
\includegraphics[width=.4\textwidth, keepaspectratio]{final_images/training_losses/lk_training_loss.png} \\
\includegraphics[width=.4\textwidth, keepaspectratio]{final_images/training_losses/mg_training_loss.png} \\
\includegraphics[width=.4\textwidth, keepaspectratio]{final_images/training_losses/time_training_loss.png} \\
\includegraphics[width=.4\textwidth, keepaspectratio]{final_images/training_losses/state_training_loss.png} \\
\includegraphics[width=.4\textwidth, keepaspectratio]{final_images/training_losses/diffusion_loss.png}
\caption{Training losses for all experiments (Lotka Volterra, Mackey-Glass, Time Dependent DDE, State Dependent DDE and Diffusion Delay PDE)}
\label{pics:training_loss}
\end{figure}

We provide the training loss of Neural Laplace where 50\% of the trajectory is fed to the model in Figure \ref{fig:50_neural_laplace_training_curve}.

\begin{figure}[H]
\centering
\includegraphics[width=.4\textwidth,keepaspectratio]{final_images/training_losses/time_training_loss_50_percent_laplace.png}
\caption{Time Dependent DDE training curve where 50\% of data is fed to Neural Laplace}
\label{fig:50_neural_laplace_training_curve}
\end{figure}


%----------------------------------------------------
\section{Model hyperparameters} 
\label{ap:model_hyper}
\label{sss:hyperparameters}

Table \ref{hyper_ivp} sums up the MLP architecture of each IVP model (ie NODE, ANODE and Neural DDE) for each dynamical systems. ANODE has an arbitrary augmented state of dimension 10 except for the PDE that has 100. Neural Laplace's architecture is the default one taken from the official implementation for all systems. The learning rate and the number of epochs are the same for all models. The optimizer used is \verb?AdaBelief? \citep{zhuang2020adabelief}. Table \ref{number_params} gives the number of parameters for each model. 
 
\begin{table}[H]
\vskip 0.15in
\begin{center}
\begin{small}
\begin{sc}
\resizebox{9cm}{!}{
\begin{tabular}{l| ccccr}

   & width & depth & activation & epochs & $l_r$  \\
\hline
Lokta-Volterra & 64 & 3 & relu & 500 & .001 \\
Mackey-Glass & 64 & 3 & relu & 1000 & .001 \\
Time Dependent DDE & 64 & 3 & relu & 2000 & .001 \\
State Dependent DDE & 64 & 3 & relu & 1000 & .001 \\
Diffusion PDE DDE & 128 & 3 & relu & 500 & .0001 \\
\hline
\end{tabular}}
\end{sc}
\end{small}
\end{center}
\caption{Model and training hyperparameters}
\label{hyper_ivp}
\end{table}

\begin{table}[H]
\vskip 0.15in
\begin{center}
\begin{small}
\begin{sc}
\resizebox{9cm}{!}{
\begin{tabular}{l|ccccr}
   & NODE & ANODE & Neural DDE & Neural Laplace  \\
\hline
Lokta-Volterra & 8642 & 9958 & 8771  & 17194 \\
Mackey-Glass & 8513 & 9815 & 8578  & 17194 \\
Time Dependent DDE & 8513 & 9815 & 8578  & 17194 \\
State Dependent DDE & 8513 & 9815 & 8642  & 17194 \\
Diffusion PDE DDE & 58852 & 84552 & 71653 & 17194 \\
\hline
\end{tabular}}
\end{sc}
\end{small}
\end{center}
\caption{Number of parameters for each DDE system}
\label{number_params}
\end{table}


%----------------------------------------------------
\section{Data generation parameters} \label{ap:data_generation}
\label{sss:dyna_recap}

We expose in Table \ref{hyper_dataset} the parameters used for each dataset generation. The start integration time is always $T_0=0$. $T_F$ refers to the end time integration. \verb?NUM_STEPS? equally spaced points are sampled in $[T_0, T_F]$. The specific delays \verb?DELAYS? and the constant history function $\phi(t)$ function domain are given. Each training dataset is comprised of $256$ datapoints and the testset of $32$ datapoints. We used our own DDE solver to generate the data (Dopri5 ODE solver \citep{dormand} was used.). We then double-checked and compared its validity with Julia’s DDE solver .$\mathcal{U}$ refers to the uniform distribution. For example, Mackey Glass’ constant history function value is uniformly sampled between 0.1 and 1.0. For the Diffusion Delay PDE the history function value in the column $\phi(t)$ refers to the constant $a$ defined in Section \ref{ssec:system_exp}.

\begin{table}[H]
\vskip 0.15in
\begin{center}
\begin{small}
\begin{sc}
\begin{tabular}{l|cccccccr}
&$T_F$ & num\_steps & delays &  $\phi(t)$\\
\hline
Lokta-Volterra & 15.0 & 200 & 0.2   & $\mathcal{U}(0.1, 2.0)$\\
Mackey-Glass & 80.0 & 400 & 10.0    & $\mathcal{U}(0.1, 1.0)$\\
Time Dependent DDE &  20.0 & 200 & $2\sin(t)$  & $\mathcal{U}(0.1, 2.0)$\\
State Dependent DDE &  10.0 & 150& $0.5 \cos(\mathbf{x}(t))$ & $\mathcal{U}(0.1, 1.0)$\\
Diffusion PDE DDE &  4.0 & 100& $1.0$ & $\mathcal{U}(0.1, 4.0)$\\
\hline
\end{tabular}
\end{sc}
\end{small}
\end{center}
\caption{Dataset generation information}
\label{hyper_dataset}
\end{table}


%----------------------------------------------------
\section{Experiment parameters for history step function} 
\label{ap:step_history_function}
\label{sss:history_function}

In table \ref{tab:exp_his}  we give the parameters used for each experiment. \textsc{\small EXTRAPOLATED} $\phi(t)$ indicates the possible value of the constant history function (first experiment).  $\tau_{\text{max}}$ and $c_0, c_1$ are described in evaluation's section \ref{ss:evaluation} (second experiment). For the Diffusion Delay PDE, as stated in Section \ref{ss:evaluation}, the other history step function is omitted. Figure \ref{fig:different_hist_mg} clearly shows how the modification of the history function can drastically change the nature of the dynamics.

\begin{table}[H]
\vskip 0.15in
\begin{center}
\begin{small}
\begin{sc}
\begin{tabular}{l|cccr}
   & extrapolated $\phi(t)$ & $\tau_{\text{max}}$ & $c_0$ &  $c_1$ \\
\hline
Lokta-Volterra & $\mathcal{U}(2.0, 3.0)$ & $0.2$ & $0.1$ & $3.0$ \\
% Dual history & $\mathcal{U}(-4.0, -2.0) \cup \mathcal{U}(2.0, 4.0)$  & $1/3$ & $-4.0$ & $4.0$ \\
Mackey-Glass &  $\mathcal{U}(1.0, 2.0)$ & $10.0$ & $0.1$ & $2.0$\\
Time Dependent DDE & $\mathcal{U}(2.0, 3.0)$ & $3.0$ & $0.1$ & $3.0$ \\
State Dependent DDE & $\mathcal{U}(-1.0, 0.1)$ & $1/2$ & $-1.0$ & $1.0$\\
\hline
\end{tabular}
\end{sc}
\end{small}
\end{center}
\caption{System specific values for each testing experiment}
\label{tab:exp_his}
\end{table}

\begin{figure}[H]
\centering
\includegraphics[width=0.5\textwidth,height=0.5\textheight,keepaspectratio]{final_images/illustration_different_history.png}
\caption{Different history functions (dotted blue line) for Mackey Glass system}
\label{fig:different_hist_mg}
\end{figure}